\documentclass{article}

% NeurIPS 2024 style (optional per professor)
\usepackage[preprint]{neurips_2024}
\usepackage[utf8]{inputenc}
\usepackage[T1]{fontenc}
\usepackage{hyperref}
\usepackage{url}
\usepackage{booktabs}
\usepackage{amsmath}
\usepackage{amssymb}
\usepackage{graphicx}
\usepackage{enumitem}

% -----------------------------------------------------------------------
\title{Modality-Conditioned Safety Representations\\
       for Vision-Language Models\\[4pt]
       \large Mid-Project Report --- CSCI 699}

\author{%
  
}

\begin{document}
\maketitle

% -----------------------------------------------------------------------
100\% AI Generated - I haven’t looked through it at all

\section{Introduction}

Vision-Language Models (VLMs) exhibit safety vulnerabilities that are qualitatively distinct from
those of their text-only counterparts. Prior work has established that safety-relevant concepts in
large language models are encoded as linear directions in activation space~\cite{nanda2023emergent,
arditi2024refusal}, enabling effective safety classification through geometric methods. However,
VLMs introduce a more complex challenge: the addition of visual inputs induces systematic shifts in
these safety representations, causing models to perceive multimodal inputs as \emph{safer} than
equivalent text-only queries~\cite{zou2025shiftdc}. This modality-induced safety distortion
underlies the effectiveness of image-based jailbreak attacks, where pairing harmful text with
visually relevant images dramatically increases attack success rates~\cite{liu2024mmsafetybench,
luo2024jailbreakv}.

Current detection methods treat safety as a single global concept, applying uniform classification
boundaries regardless of whether the threat originates from text, images, or their interaction. We
argue that effective VLM safety detection requires \textbf{modality-conditioned safety
representations}---boundaries that adapt based on the dominant modality contributing to a given
input. This approach aims to improve both sensitivity (catching image-based jailbreaks that global
methods miss) and specificity (reducing over-refusal of visually alarming but contextually
legitimate queries).

% -----------------------------------------------------------------------
\section{Related Work}

\paragraph{Linear Representation Hypothesis.}
Nanda et al.~\cite{nanda2023emergent} and Arditi et al.~\cite{arditi2024refusal} demonstrated that
high-level semantic concepts---including safety features---are encoded as linear directions in
transformer activation spaces. Crucially, Arditi et al.\ showed that refusal behavior in LLMs is
mediated by a \emph{single} linear direction, and that ablating this direction reliably suppresses
refusal across diverse harmful inputs. This geometric structure underpins our representation-based
approach.

\paragraph{Modality-Induced Safety Distortion.}
ShiftDC~\cite{zou2025shiftdc} demonstrates that visual inputs cause consistent activation shifts
that distort safety perception: multimodal inputs appear safer than text-only equivalents in
hidden-state geometry even when semantic content is held fixed. The authors propose a global
projection-based correction; a key limitation is that this correction is modality-agnostic. Our
work addresses this by learning \emph{separate} classification boundaries per modality regime.

\paragraph{Geometric Safety Detection.}
RCS~\cite{liu2025rcs} frames jailbreak detection as contrastive boundary detection in activation
space. While effective for text-dominant threats, it does not condition on modality. SPIN~\cite
{chen2025spin} identifies that attention imbalance---text tokens attend to image tokens less than
10\% of the time despite images comprising 76--92\% of input---contributes to safety
vulnerabilities, motivating our use of attention allocation as a conditioning signal.

\paragraph{Benchmarks.}
HoliSafe-Bench~\cite{lee2025holisafe} provides per-modality safety labels across five image-text
combinations (e.g., unsafe image + safe text), making it uniquely suited for our differential
analysis. MSSBench~\cite{zhou2024mss} evaluates situational safety; MM-SafetyBench~\cite
{liu2024mmsafetybench} and JailBreakV-28K~\cite{luo2024jailbreakv} target adversarial attacks.

% -----------------------------------------------------------------------
\section{Technical Approach}

Our goal is to learn modality-conditioned safety subspaces $\{\mathcal{S}_m\}$ such that, given an
input $\mathbf{x} = (\mathbf{x}_\text{img}, \mathbf{x}_\text{txt})$, safety classification uses a
boundary adapted to the dominant threat modality.

\subsection{Step 1: Characterizing Modality-Specific Safety Directions}
We extract intermediate-layer activations $\mathbf{h}^{(\ell)}\in\mathbb{R}^d$ from a frozen VLM
for inputs belonging to each of the five HoliSafe-Bench modality combinations. For combination
$(i,j)$, we compute a safety direction:
\begin{equation}
  \mathbf{r}_{ij} = \mu_\text{unsafe}^{(ij)} - \mu_\text{safe}^{(ij)},
\end{equation}
where $\mu$ is the mean activation over the respective subset. We then analyze pairwise cosine
similarities $\langle\hat{\mathbf{r}}_{ij},\hat{\mathbf{r}}_{i'j'}\rangle$ and use PCA to
visualize whether modality conditions induce distinguishable geometric structure.

\subsection{Step 2: Modality Conditioning via Attention Signals}
Following SPIN~\cite{chen2025spin}, we use attention allocation as a lightweight conditioning
signal. Let $\alpha_\text{img}^{(\ell)}$ be the fraction of attention mass on image tokens in
layer $\ell$. We define:
\begin{equation}
  s_\text{mod}(\mathbf{x}) = \frac{1}{L}\sum_{\ell=1}^{L}\alpha_\text{img}^{(\ell)},
\end{equation}
and interpolate between two learned projection matrices:
\begin{equation}
  \mathbf{P}(\mathbf{x}) = (1-s_\text{mod})\,\mathbf{P}_\text{txt}
                          + s_\text{mod}\,\mathbf{P}_\text{img}.
\end{equation}
Safety is then classified via
$f(\mathbf{x}) = \operatorname{sign}\!\bigl(\mathbf{P}(\mathbf{x})\,\mathbf{h}^{(\ell^*)}\cdot\mathbf{w}+b\bigr)$,
where $\ell^*$ is selected by validation performance.

\subsection{Step 3: Learning Objective}
We jointly learn $\mathbf{P}_\text{txt}$ and $\mathbf{P}_\text{img}$ to maximize inter-class
separation within each modality regime:
\begin{equation}
  \mathcal{L} = \mathcal{L}_\text{contrastive}
              + \lambda\,\|\mathbf{P}_\text{txt}-\mathbf{P}_\text{img}\|_F,
\end{equation}
where the regularization term encourages the two projections to differ meaningfully.

\subsection{Implementation Details}
\begin{itemize}[leftmargin=*]
  \item \textbf{Base VLM:} % TODO: e.g., LLaVA-1.5-7B
  \item \textbf{Primary benchmark:} HoliSafe-Bench (5 modality splits)
  \item \textbf{Secondary:} MSSBench, MM-SafetyBench
  \item \textbf{Metrics:} Attack Success Rate (ASR $\downarrow$) on unsafe splits;
        Refusal Rate (RR $\downarrow$) on safe splits
  \item \textbf{Hardware:} % TODO: GPU setup
\end{itemize}

% -----------------------------------------------------------------------
\section{Progress and Preliminary Results}

\subsection{Experimental Setup}
% TODO: dataset access, model checkpoint, extraction procedure, framework used

\subsection{Preliminary Observations}
% TODO: key findings so far, e.g.:
% - PCA plots of activations across modality conditions
% - Cosine similarities between safety directions r_ij
% - Baseline ASR of undefended model on each HoliSafe-Bench split
\textbf{[To be completed with experimental results.]}

\subsection{Baseline Comparison}

\begin{table}[h]
\centering
\caption{Preliminary results on HoliSafe-Bench (\%). Dashes = not yet available.}
\label{tab:results}
\begin{tabular}{lcccc}
\toprule
 & \multicolumn{2}{c}{Unsafe Img + Safe Text}
 & \multicolumn{2}{c}{Safe Img + Unsafe Text} \\
\cmidrule(lr){2-3}\cmidrule(lr){4-5}
Method & ASR$\downarrow$ & RR$\downarrow$ & ASR$\downarrow$ & RR$\downarrow$ \\
\midrule
No defense                           & --- & --- & --- & --- \\
ShiftDC~\cite{zou2025shiftdc}        & --- & --- & --- & --- \\
RCS~\cite{liu2025rcs}               & --- & --- & --- & --- \\
\textbf{Ours}                        & --- & --- & --- & --- \\
\bottomrule
\end{tabular}
\end{table}

% -----------------------------------------------------------------------
\section{Challenges and Risks}

\begin{enumerate}[leftmargin=*]
  \item \textbf{Benchmark access.}
        % TODO: describe any issues with dataset access or preprocessing.

  \item \textbf{Layer selection sensitivity.}
        Safety directions vary substantially across transformer layers. Selecting $\ell^*$ without
        overfitting to validation data requires careful cross-layer analysis.

  \item \textbf{Generalization of linear interpolation.}
        The interpolation scheme assumes a smooth modality spectrum, which may not hold for
        compositional attacks (safe image + safe text $\to$ unsafe combined). We may need
        non-linear conditioning or a discrete three-way classifier.

  \item \textbf{Compute constraints.}
        Activation extraction on a 7B-parameter model across large benchmarks requires substantial
        GPU memory.
        % TODO: describe current setup and mitigation strategy.

  \item \textbf{Scope management.}
        The project spans direction analysis, projection learning, and conditioning signal design.
        We will prioritize completing Step 1 first to ensure clean findings regardless of timeline.
\end{enumerate}

% -----------------------------------------------------------------------
\section{Next Steps and Timeline}

\begin{table}[h]
\centering
\caption{Remaining project milestones.}
\label{tab:timeline}
\begin{tabular}{lp{9.5cm}}
\toprule
Dates & Milestone \\
\midrule
Mar 31 -- Apr 6
  & Complete activation extraction on HoliSafe-Bench; compute and compare per-modality
    safety directions $\mathbf{r}_{ij}$ across layers. \\
Apr 7 -- Apr 13
  & Implement modality-conditioned projection learning; establish baseline comparisons
    with ShiftDC and RCS. \\
Apr 14 -- Apr 20
  & Evaluate on MSSBench and MM-SafetyBench; run ablation studies (layer choice,
    $\lambda$, linear vs.\ nonlinear conditioning). \\
Apr 21 -- Apr 27
  & PCA visualizations, error analysis, and draft final report. \\
Apr 28 -- May 4
  & Final report revisions and submission. \\
\bottomrule
\end{tabular}
\end{table}

% -----------------------------------------------------------------------
\bibliographystyle{plain}
\bibliography{references}

\end{document}
